% !TeX root = ../easyfloats.tex

\section{Other packages}
\label{other-packages}

This section lists other packages which may be useful in combination with this package.

\pkgdoc{placeins}
  When loaded with the `section` package option it prevents floats from floating to another section.
  It provides the `\FloatBarrier` command which prevents floats from floating past a certain point.
\pkgdoc{flafter}
  ensures that floats are not placed before their inclusion in the source code.
  (With the placement=t it is possible that they are placed on the top of the same page.)
\pkgdoc{booktabs}
  for formatting tables.
  This is loaded by default unless you use the package option \pkgoptn{nobooktabs}.
  Please read sections 1 and~2 of the \pkg{booktabs} documentation~\autocite{booktabs}.
\pkgdoc{xcolor}
  When loaded with the package option \pkgoptn{table} it provides commands for coloring tables.

  `\rowcolor{<color>}` sets a background color for a single row. See \cref{local-definitions-in-tables}.

  `\rowcolors{<firstrow>}{<oddcolor>}{<evencolor>}` can be used with the `exec` key and sets alternating row colors for the entire table.

\pkgdoc{array}
  extends the column specification syntax and defines the `\newcolumntype` command to define custom column types.
  It also changes the approach how rules are drawn but that is irrelevant if you apply \pkg{booktabs}' guidelines~\autocite[section~2 \sectionname{The layout of formal tables}]{booktabs}.
  This is loaded by default unless you use the package option \pkgoptn{noarray}.
\pkgdoc{siunitx}
  for typesetting numbers and units.
  It provides the `S` column to align numbers at their decimal separator.
\pkgdoc{tabularx}
  A table where the columns adapt to the width of the table, not the other way around.
  Unlike `tabular*` the space goes into the columns, not between the columns.
\pkgdoc{longtable}
  provides tables where a pagebreak is allowed,
  see \cref{longtable}.
  This can be loaded automatically with the package option \pkgoptn{longtable}.
\pkgdoc{hyperref}
  automatically creates links in the pdf document for example from references to floating objects.
  By default links are set in a red box, you can disable the red box with the package option `hidelinks`.
  With the package option `pdfusetitle` it automatically sets the pdf title and author based on `\title` and `\author`.
  \pkg{hyperref} should usually be loaded last.
\pkgdoc{varioref}
  automatically inserts a page reference (e.g. \enquote{on the following page} or \enquote{on page~42}) after the reference if the reference is on a different page than the label when using `\vref{<label>}`/`\Vref{<label>}`, see the example in \cref{vref}.
\pkgdoc{cleveref}
  automatically inserts the type of reference (e.g. \enquote{figure} or \enquote{section}) in front of the reference when using `\cref{<label>}`/`\Cref{<label>}` (defined by \pkg{cleveref}) or `\vref{<label>}`/`\Vref{<label>}` (defined by \pkg{varioref}).
  With the `nameinlink` package option the type in front of the number becomes part of the link created by hyperref.
  By default `\cref`/`\vref` (to be used inside of a sentence) use abbreviations but `\Cref`/`\Vref` (to be used at the beginning of a sentence) do not because abbreviating at the beginning of a sentence is considered bad style~\autocite{texexchange_abbrev_bad_style}.
  Abbreviations can be disabled with the `noabbrev` option.
  \pkg{cleveref} must be loaded last, even after \pkg{hyperref}.
  See the example in \cref{vref}.
\pkgdoc{biblatex}
  If you input graphics you need to specify the source.
  Biblatex creates an entire bibliography for you.
\pkgdoc{tikz}
  is an amazingly powerful package to create your own graphics in \LaTeX.
\pkgdoc{newfloat} provides a more modern command to define new floating environments than the \pkg{float} package.
\pkgdoc{minted} provides syntax highlighting for source code.
\endpkgdoc


For more information about floats see \url{https://latexref.xyz/Floats.html} (it seems this is an html version of the above quoted pdf~\autocite{latex2e}).
